\documentclass{homework}

% Put in your Name !
\newcommand{\hwname}{Stella Srzich, 3371930}

\usepackage{color,soul}
\usepackage{amsmath}
\usepackage{enumitem}
\allowdisplaybreaks

% Header
\newcommand{\hwtype}{Assignment}
\newcommand{\hwclass}{Basin MCMC}
\newcommand{\hwnum}{1}
\newcommand{\duedate}{Friday, 11 September 2026 4PM}

%%% Define shortcuts here
\newcommand{\B}{\mathcal B}
% instead of typing \mathcal B to display a curly B I can now just type \B
\newcommand{\ip}[2]{\left\langle #1\middle| #2 \right\rangle}
% this defines the new command \ip{x}{y} to display <x,y>

\newcommand*{\N}{\mathbb N}
\newcommand*{\Z}{\mathbb Z}
\newcommand*{\Q}{\mathbb Q}
\newcommand*{\R}{\mathbb R}
\newcommand*{\C}{\mathbb C}
\newcommand{\F}{\mathbb F}
\newcommand{\norm}[1]{\|  #1 \|}
\newcommand{\abs}[1]{\left|#1\right|}

\newcommand{\sech}{\mathrm{sech}}
\newcommand{\csch}{\mathrm{csch}}

% commands I have added
\newcommand*{\e}{\mathrm{e}}
\newcommand*{\im}{\mathrm{i}}
\newcommand*{\cint}{\oint\limits}
\newcommand*{\rmd}{\mathrm{d}}
\newcommand*{\D}{\mathcal{D}}
\newcommand{\Pm}{\mathbb{P}} % for primes

%%%%
\begin{document}

\maketitle

The purpose of this assignment is to break down the 
development of the initial basin probability distribution and probability 
in the Basin MCMC algorithm.

Throughout this assignment, we use the following notation:
\[
\mathcal{X} := \bigcup_{n=0}^\infty \mathcal{X}_n,
\]
where
\[
    \mathcal{X}_n := \left( [0, X_{\text{max}} ] \times (0, D_{\text{max}} ] \times [0, \infty) \times [-1, 1] \right)^n.
\]


% -------- QUESTION 1 ---------
\question*{}  
Suppose we have a representative atom $\theta_0=(x_{00},d_0,R_0,a_0)^{\mathsf T}\in\mathcal{X}_1$.
Provide a function that approximates the log posterior $\log p(\theta\mid y)$ on $\mathcal{X}_1$, 
where the observed data is the noise-free data $y=A(\theta_0)$ generated by this atom.

The function $f_1: \mathcal{X}_1 \to \R$ should satisfy the following properties:
\begin{enumerate}
    \item unimodal;
    \item a good approximation to the log posterior in the region of high probability around $\theta_0$;
    \item easy to compute defining quantities from $\theta_0$;
    \item easy to sample from;
    \item as simple as possible.
\end{enumerate}

Note that this function need only be specified up to an additive constant.

\textit{Solution.} 
We present the solution as follows:
\begin{enumerate}
    \item Define $f_1$ as the second-order Taylor polynomial of the log posterior about $\theta_0$;
    \item Derive a formula for the gradient $g=\nabla\log p(\theta_0\mid y)$;
    \item Derive a formula for the negative Hessian $P=-\nabla^2\log p(\theta_0\mid y)$;
    \item State the final formula for $f_1$.
\end{enumerate}
\begin{enumerate}
    \item 
    For $\theta \in \mathcal{X}_1$, the log posterior is given by 
    \[
    \log p(\theta\mid y)
    =
    -\frac{1}{2\sigma^2}\lVert y-A(\theta)\rVert_2^2
    +\log p(\theta) + C,
    \]
    where $C$ is some constant independent of $\theta$, 
    and $A: \mathcal{X} \to [0, 1]^{N_t N_x}$ is the forward map.
    Performing a second-order Taylor expansion about $\theta_0$ 
    (since $\log p(\theta\mid y) \in C^3(U)$ for some open neighbourhood $U\subset\mathcal{X}_1$ of $\theta_0$, provided $\theta_0\in\text{int}(\mathcal{X}_1)$) gives
    \begin{align}
    \log p(\theta\mid y)
    = 
    \log p(\theta_0\mid y)
    +g^{\mathsf T}\left(\theta-\theta_0\right)
    -\frac12\left(\theta-\theta_0\right)^{\mathsf T}P\left(\theta-\theta_0\right)
    + O(\norm{\theta-\theta_0}^3), \label{eq:taylor_expansion_simple}
    \end{align}
    where
    \[
    g:=\nabla\log p(\theta_0\mid y) \quad \text{and} \quad P:=-\nabla^2\log p(\theta_0\mid y).
    \]
    Motivated by \eqref{eq:taylor_expansion_simple}, we define our approximation to the log posterior $f_1: \mathcal{X}_1 \to \R$ by
    \[
    f_1(\theta) := \log p(\theta_0\mid y)
    +g^{\mathsf T}\left(\theta-\theta_0\right)
    -\frac12\left(\theta-\theta_0\right)^{\mathsf T}P\left(\theta-\theta_0\right).
    \]

    \item We have that
    \begin{align}
    g 
    = \nabla\log p(\theta_0\mid y)
    = \nabla\left( \log p(y\mid\theta_0) + \log p(\theta_0) - \log p(y) \right)
    =
    \nabla\log p(y\mid\theta_0)
    +
    \nabla\log p(\theta_0). \label{eq:gradient_log_posterior}
    \end{align}

    For $\theta \in \mathcal{X}_1$, we have that
    \begin{align*}
    \nabla\log p(y\mid\theta)
    = \nabla\left(-\frac{1}{2\sigma^2}\lVert y-A(\theta)\rVert_2^2 \right)
    = \frac{1}{\sigma^2}J_A(\theta)^{\mathsf T}\left(y-A(\theta) \right),
    \end{align*}
    where $J_A(\theta)$ 
    is the Jacobian of the forward map $A$ (restricted to the domain $\mathcal{X}_1$).
    Recalling that $y=A(\theta_0)$,
    \begin{align}
    \nabla\log p(y\mid\theta_0) = \frac{1}{\sigma^2}J_A(\theta_0)^{\mathsf T}\left(y-A(\theta_0) \right) = 0. \label{eq:gradient_log_likelihood_0}
    \end{align}

    Finally, for $\theta \in \mathcal{X}_1$, we have that
    \begin{align}
    \nabla \log p(\theta) 
    = \nabla \log \left(\lambda_R \e^{-\lambda_r R} \right)
    = \nabla \left(-\lambda_R R \right)
    = \left(0, 0, -\lambda_R, 0 \right)^{\mathsf T}. \label{eq:gradient_log_prior}
    \end{align}
    
    Combining \eqref{eq:gradient_log_posterior}, \eqref{eq:gradient_log_likelihood_0}, and \eqref{eq:gradient_log_prior}, we have that
    \[
    g
    =
    \left(0, 0, -\lambda_R, 0 \right)^{\mathsf T}.
    \]

    \item The negative Hessian of the log posterior is given by
    \begin{align}
    P = -\nabla^2\log p(\theta_0\mid y)
    =
    -\nabla^2\log p(y\mid\theta_0)
    -\nabla^2\log p(\theta_0). \label{eq:hessian_log_posterior}
    \end{align}

    Recalling \eqref{eq:gradient_log_prior}, we have that
    \begin{align}
    \nabla^2\log p(\theta_0)=0. \label{eq:hessian_log_prior_0}
    \end{align}

    Next, for $\theta \in \mathcal{X}_1$, one can show that
    \begin{align*}
    \nabla^2\log p(y\mid\theta)
    =
    -\frac{1}{\sigma^2}J_A(\theta)^{\mathsf T}J_A(\theta)
    +
    \frac{1}{\sigma^2}
    \sum_k (y - A(\theta))_k \nabla^2A_k(\theta).
    \end{align*}
    Again recalling that $y = A(\theta_0)$, we obtain
    \begin{align}
    \nabla^2\log p(y\mid\theta_0)
    = -\frac{1}{\sigma^2}J_A(\theta_0)^{\mathsf T}J_A(\theta_0). \label{eq:hessian_log_likelihood_0}
    \end{align}

    Combining \eqref{eq:hessian_log_posterior}, \eqref{eq:hessian_log_prior_0}, and \eqref{eq:hessian_log_likelihood_0}, we have that
    \begin{align}
    P
    =
    \frac{1}{\sigma^2}J_A(\theta_0)^{\mathsf T}J_A(\theta_0). \label{eq:hessian_log_posterior_0}
    \end{align}

    To obtain a formula for \(P\), we will obtain a formula for $J_A(\theta_0)$. 
    To this end, recall that the $t N_x + x$-th entry of the forward map $A: \mathcal{X} \to [0, 1]^{N_t N_x}$ 
    with parameters $\theta=(x_0, d, R, a)^{\mathsf T} \in \mathcal{X}_1$ is given by 
    \begin{align*}
    A(t, x; \theta) = aB(x)w(t-T(x))
    \end{align*}
    where
    \begin{align*}
    B(x) = \sqrt{\frac{d+R}{s(x)}} e^{-2\alpha s(x)}, \quad T(x) = \frac{2}{v} (s(x) - (d+R)) + \frac{2d}{v} = \frac{2}{v} (s(x) - R)
    \end{align*}
    with $s(x) = \sqrt{(d+R)^2+(x-x_0)^2}$,
    and 
    \begin{align*}
    w(\tau) = (1-2(\pi f_0\tau)^2)e^{-(\pi f_0\tau)^2}.
    \end{align*}

    Applying the product and chain rules, we obtain
    \begin{equation}
    \begin{aligned}
    \frac{\partial A(t,x;\theta)}{\partial x_0}&=aB(x)\left[w(t-T(x))\frac{\partial\log B(x)}{\partial x_0}+w'(t-T(x))\frac{-\partial T(x)}{\partial x_0}\right],\\
    \frac{\partial A(t,x;\theta)}{\partial d}&=aB(x)\left[w(t-T(x))\frac{\partial\log B(x)}{\partial (d+R)}+w'(t-T(x))\frac{-\partial T(x)}{\partial d}\right],\\
    \frac{\partial A(t,x;\theta)}{\partial R}&=aB(x)\left[w(t-T(x))\frac{\partial\log B(x)}{\partial (d+R)}+w'(t-T(x))\frac{-\partial T(x)}{\partial R}\right],\\
    \frac{\partial A(t,x;\theta)}{\partial a}&=B(x)\,w(t-T(x)).
    \end{aligned}
    \label{eq:forward_map_partials}
    \end{equation}
    
    The needed derivatives are given by
    \begin{align*}
    w'(\tau)&=2(\pi f_0)^2\tau(2(\pi f_0\tau)^2-3)e^{-(\pi f_0\tau)^2},\\
    \frac{\partial\log B(x)}{\partial x_0}&=(x-x_0)\left(\frac{1}{2s(x)^2}+\frac{2\alpha}{s(x)}\right),\\
    \frac{\partial\log B(x)}{\partial (d+R)}&=\frac{1}{2(d+R)}-\frac{(d+R)}{2s(x)^2}-\frac{2\alpha (d+R)}{s(x)},\\
    -\frac{\partial T(x)}{\partial x_0}&=\frac{2(x-x_0)}{vs(x)},\\
    -\frac{\partial T(x)}{\partial d}&=-\frac{2(d+R)}{v s(x)},\\
    -\frac{\partial T(x)}{\partial R}&=\frac{2}{v}\left(1-\frac{(d+R)}{s(x)}\right).
    \end{align*}
    
    Vectorising each of the derivatives in \eqref{eq:forward_map_partials} 
    gives the four columns of \(J_A(\theta_0)\).

    \item Recall that for $\theta \in \mathcal{X}_1$, our approximation to the log posterior is given by
    \begin{align}
    f_1(\theta) 
    &= \log p(\theta_0\mid y)+g^{\mathsf T}\left(\theta-\theta_0\right)-\frac12\left(\theta-\theta_0\right)^{\mathsf T}P\left(\theta-\theta_0\right).\nonumber
    \end{align}
    Assuming that $P$ is positive definite and so $P^{-1}$ exists (which holds if and only if $J_A(\theta_0)$ has full rank), we may complete the square to obtain
    \begin{align}
    f_1(\theta) 
    &= \log p(\theta_0\mid y)+g^{\mathsf T}\left(\theta-\theta_0\right)-\frac12\left(\theta-\theta_0\right)^{\mathsf T}P\left(\theta-\theta_0\right)\nonumber\\
    &= \log p(\theta_0\mid y)-\frac12\left[(\theta-\theta_0)^\top P(\theta-\theta_0)-2g^\top(\theta-\theta_0)\right]\nonumber\\
    &= \log p(\theta_0\mid y)-\frac12\left[(\theta-\theta_0)^\top P(\theta-\theta_0)-2g^\top(\theta-\theta_0)+g^\top P^{-1}g-g^\top P^{-1}g\right]\nonumber\\
    &= \log p(\theta_0\mid y)-\frac12\left[(\theta-\theta_0-P^{-1}g)^\top P(\theta-\theta_0-P^{-1}g)-g^\top P^{-1}g\right]\nonumber\\
    &= \log p(\theta_0\mid y) + \frac12g^\top P^{-1}g - \frac12(\theta-\theta_0-P^{-1}g)^\top P(\theta-\theta_0-P^{-1}g)\nonumber\\
    &= C' - \frac12(\theta-\mu)^\top P(\theta-\mu), \nonumber
    \end{align}
    where $C'$ is some constant independent of $\theta$, and $\mu := \theta_0+P^{-1}g$.

    Thus, our approximation $f_1(\theta)$ defines a truncated Gaussian on $\mathcal{X}_1$, 
    with mean $\mu$ and covariance $P^{-1}$.
    In the implementation, the Moore--Penrose pseudoinverse \(P^{+}\) is used in
    place of \(P^{-1}\) for numerical robustness.
\end{enumerate}


% -------- QUESTION 2 ---------
\pagebreak
\question*{}  
Suppose we have a representative atom $\theta_0=(x_{00},d_0,R_0,a_0)^{\mathsf T}\in\mathcal{X}_1$.
Provide a function that approximates the log likelihood $\log p(y\mid\theta)$ on $\mathcal{X}_2$, 
where the observed data is the noise-free data $y=A(\theta_0)$ generated by this atom.

The function $\tilde{f}_2: \mathcal{X}_2 \to \R$ should satisfy the following properties:
\begin{enumerate}
    \item unimodal;
    \item exchangeable, i.e., $\tilde{f}_2(\theta_1, \theta_2) = \tilde{f}_2(\theta_2, \theta_1)$;
    \item a good approximation to the log likelihood in the region of high probability around $\theta_0$;
    \item easy to compute defining quantities from $\theta_0$;
    \item easy to sample from;
    \item as simple as possible.
\end{enumerate}

Note that this function need only be specified up to an additive constant.

\textit{Solution.} 
We present the solution as follows: \hl{TODO: Change the presentation once the solution is complete.}
\begin{enumerate}
    \item Present the log likelihood and the forward map;
    \item Define a reparameterization;
    \item Rewrite the forward map in terms of the reparameterization;
    \item Taylor approximate the forward map around the center of the reparameterization;
    \item Taylor approximate the forward map around the representative atom;
    \item Approximate the Hessian of the forward map around the center of the reparameterization by the Hessian of the forward map around the representative atom;
    \item Combine the Taylor approximations to obtain the final formula for $\tilde{f}_2$.
\end{enumerate}
\begin{enumerate}
    \item For $(\theta_1, \theta_2) \in \mathcal{X}_2$, the log likelihood is given by 
    \begin{align*}
    \log p(y\mid\theta_1, \theta_2)
    = - \frac{1}{2\sigma^2} \norm{A(\theta_0) - A(\theta_1) - A(\theta_2)}_2^2 + C
    \end{align*}
    where $C$ is some constant independent of $(\theta_1, \theta_2)$, and $A: \mathcal{X} \to [0, 1]^{N_t N_x}$ is the forward map.
    Recall that the $t N_x + x$-th entry of the forward map $A: \mathcal{X} \to [0, 1]^{N_t N_x}$ 
    with parameters $\theta=(x_0, d, R, a)^{\mathsf T} \in \mathcal{X}_1$ is given by 
    \begin{align*}
    A(t, x; \theta) = aB(x)w(t-T(x))
    \end{align*}
    where
    \begin{align*}
    B(x) = \sqrt{\frac{d+R}{s(x)}} e^{-2\alpha s(x)}, \quad T(x) = \frac{2}{v} (s(x) - (d+R)) + \frac{2d}{v} = \frac{2}{v} (s(x) - R)
    \end{align*}
    with $s(x) = \sqrt{(d+R)^2+(x-x_0)^2}$,
    and 
    \begin{align*}
    w(\tau) = (1-2(\pi f_0\tau)^2)e^{-(\pi f_0\tau)^2}.
    \end{align*}

    \item Define $g_0:=(x_{00},d_0,R_0)^{\mathsf T}$, so that we may write the representative atom as 
    $\theta_0=(g_0,a_0)^{\mathsf T}$.
    Now, with two atoms $\theta_1=(g_1,a_1)^{\mathsf T}, \theta_2=(g_2,a_2)^{\mathsf T} \in \mathcal{X}_1$, 
    where $g_1:=(x_1,d_1,R_1)^{\mathsf T}$ and $g_2:=(x_2,d_2,R_2)^{\mathsf T}$, 
    and $a_1 + a_2 \neq 0$,
    we will use the following reparameterization:
    \begin{itemize}
    \item $a := a_1 + a_2$ is the total amplitude,
    \item $\lambda := \frac{a_1}{a_1 + a_2}$ is the amplitude split,
    \item $g := \lambda g_1 + (1-\lambda)g_2$ is the amplitude-weighted center,
    \item $h := g_1 - g_2$ is the difference.
    \end{itemize}
    We may recover the original parameters as follows:
    \begin{itemize}
    \item $a_1 = \lambda a$,
    \item $a_2 = (1-\lambda) a$,
    \item $g_1 = g + (1-\lambda)h$,
    \item $g_2 = g - \lambda h$.
    \end{itemize}

    \item Define the function $F: [0,X_{\text{max}}]\times (0,D_{\text{max}}] \times[0, \infty) \to [0,1]^{N_t N_x}$ by 
    \[F(g) := A(g, 1),\] 
    where $A(g,1)$ is the $N_t N_x$-dimensional vector with $A(t,x;g,1)$ as the $t N_x + x$-th entry.
    Notice that $A(g, a) = aF(g)$ for $a \in [-1, 1]$. Hence, 
    \begin{align}
    A(\theta_1) + A(\theta_2) 
    &= a_1F(g_1) + a_2F(g_2) \nonumber\\
    &= a\left[\lambda F(g_1) + (1-\lambda )F(g_2)\right] \nonumber\\
    &= a\left[\lambda F(g + \delta_1) + (1-\lambda )F(g + \delta_2)\right], \label{eq:sum_forward_map}
    \end{align}
    where $\delta_1 := g_1 - g = (1-\lambda)h$ and $\delta_2 := g_2 - g = -\lambda h$.

    \item Now, for $g \in \operatorname{int}\left({\mathcal{X}_1}\right)$, 
    (since $F$ is smooth in an open neighbourhood of $g\in\operatorname{int}\left({\mathcal{X}_1}\right)$) 
    performing a second-order Taylor expansion of $F$ around $g$ gives
    \begin{align}
    F(g + \delta)
    &= F(g) + J_F(g) \delta + \frac{1}{2} H_F(g)[\delta, \delta] + O(\norm{\delta}^3) \label{eq:taylor_expansion}
    \end{align}
    as $\delta \to 0$, where $J_F(g)$ is the Jacobian of $F$ at $g$,
    \[
    J_F(g) = \left[
    \frac{\partial F(g)}{\partial x_0},
    \frac{\partial F(g)}{\partial d},
    \frac{\partial F(g)}{\partial R}
    \right]
    =
    \begin{bmatrix}
    \frac{\partial F_1(g)}{\partial x_0}
    &
    \frac{\partial F_1(g)}{\partial d}
    &
    \frac{\partial F_1(g)}{\partial R}
    \\
    \vdots & \vdots & \vdots
    \\
    \frac{\partial F_{N_tN_x}(g)}{\partial x_0}
    &
    \frac{\partial F_{N_tN_x}(g)}{\partial d}
    &
    \frac{\partial F_{N_tN_x}(g)}{\partial R}
    \end{bmatrix},
    \]
    and $H_F(g)[\delta, \delta]$ is the $N_t N_x$-dimensional vector with $k$-th entry
    \[
    \left(H_F(g)[\delta, \delta]\right)_k = \delta^\mathsf T \nabla^2 F_k(g) \delta,
    \]
    where $F_k(g)$ is the $k$-th entry of $F(g)$.

    Using the Taylor approximation \eqref{eq:taylor_expansion} for each term in \eqref{eq:sum_forward_map} gives
    \begin{align}
    \lambda F(g + \delta_1) + (1-\lambda)F(g + \delta_2)
    &= F(g) + J_F(g) \left(\lambda\delta_1 + (1-\lambda)\delta_2\right) + \frac{1}{2} \left(\lambda H_F(g)[\delta_1, \delta_1]+ (1-\lambda) H_F(g)[\delta_2, \delta_2]\right) \nonumber\\
    &+ O(\norm{\delta_1}^3 + \norm{\delta_2}^3), \label{eq:sum_forward_map_taylor}
    \end{align}
    as $\left(\delta_1, \delta_2\right) \to 0$.
    Now, notice that since $\lambda\delta_1 + (1-\lambda)\delta_2 = 0$, we have that
    \begin{align}
    J_F(g) \left(\lambda \delta_1 + (1-\lambda)\delta_2\right) = 0. \label{eq:jacobian_zero}
    \end{align}
    Next, notice that since $H_F(g)[\delta, \delta]$ is bilinear, we have that
    \begin{align}
    \lambda H_F(g)[\delta_1, \delta_1]+ (1-\lambda) H_F(g)[\delta_2, \delta_2]
    &= \lambda (1-\lambda)^2 H_F(g)[h, h]+ (1-\lambda) \lambda^2 H_F(g)[h, h] \nonumber\\
    &= \lambda(1-\lambda)H_F(g)[h, h]. \label{eq:bilinear_property}
    \end{align}
    Combining \eqref{eq:sum_forward_map_taylor}, \eqref{eq:jacobian_zero}, and \eqref{eq:bilinear_property} gives
    \begin{align}
    \lambda F(g + \delta_1) + (1-\lambda)F(g + \delta_2)
    &= F(g) + \frac{1}{2} \lambda(1-\lambda)H_F(g)[h, h] + O(\norm{\delta_1}^3 + \norm{\delta_2}^3). \label{eq:sum_forward_map_taylor_final}
    \end{align}
    Since
    \[
    \norm{\delta_1}^3 + \norm{\delta_2}^3 = \left(|\lambda|^3 + |1-\lambda|^3\right)\norm{h}^3,
    \]
    if we restrict to the region where $|a_1+a_2|>a_\text{min}$ for some $a_\text{min} > 0$ 
    \hl{(which should be reasonable otherwise they cancel each other out in the image)}, 
    so that $|\lambda|, |1-\lambda| < 1/a_\text{min}$, then we have that 
    \begin{align}
    O(\norm{\delta_1}^3 + \norm{\delta_2}^3) = O\left(\norm{h}^3\right), \label{eq:error_bound}
    \end{align}
    as $h \to 0$.
    Overall, combining \eqref{eq:sum_forward_map_taylor_final} and \eqref{eq:error_bound} gives
    \begin{align}
    \lambda F(g + \delta_1) + (1-\lambda)F(g + \delta_2)
    &= F(g) + \frac{1}{2} \lambda(1-\lambda)H_F(g)[h, h] + O\left(\norm{h}^3\right) \label{eq:sum_forward_map_taylor_final_final}
    \end{align}
    as $h \to 0$ (where the $O\left(\norm{h}^3\right)$ remainder is uniform in $\lambda$ in the region where $|a_1+a_2|>a_\text{min}$).

    \item 
    Now, for $g_0 \in \operatorname{int}\left({\mathcal{X}_1}\right)$, 
    (since $F$ is smooth in an open neighbourhood of $g_0\in\operatorname{int}\left({\mathcal{X}_1}\right)$) 
    performing a first-order Taylor expansion of $F$ around $g_0$ gives
    \begin{align}
    F(g_0 + \delta)
    &= F(g_0) + J_F(g_0) \delta + O(\norm{\delta}^2) \label{eq:taylor_expansion_2}
    \end{align}
    as $\delta \to 0$, where $J_F(g_0)$ is the Jacobian of $F$ at $g_0$, given by
    \[
    J_F(g_0) = \left[
    \frac{\partial F(g_0)}{\partial x_0},
    \frac{\partial F(g_0)}{\partial d},
    \frac{\partial F(g_0)}{\partial R}
    \right]
    =
    \begin{bmatrix}
    \frac{\partial F_1(g_0)}{\partial x_0}
    &
    \frac{\partial F_1(g_0)}{\partial d}
    &
    \frac{\partial F_1(g_0)}{\partial R}
    \\
    \vdots & \vdots & \vdots
    \\
    \frac{\partial F_{N_tN_x}(g_0)}{\partial x_0}
    &
    \frac{\partial F_{N_tN_x}(g_0)}{\partial d}
    &
    \frac{\partial F_{N_tN_x}(g_0)}{\partial R}
    \end{bmatrix}.
    \]
    Applying the approximation \eqref{eq:taylor_expansion_2} to $F(g)$ in \eqref{eq:sum_forward_map_taylor_final_final} gives
    \begin{align}
    \lambda F(g + \delta_1) + (1-\lambda)F(g + \delta_2)
    &= F(g_0 + \delta_0) + \frac{1}{2} \lambda(1-\lambda)H_F(g)[h, h] + O\left(\norm{h}^3\right) \nonumber\\
    &= F(g_0) + J_F(g_0) q + \frac{1}{2} \lambda(1-\lambda)H_F(g)[h, h] + O\left(\norm{\delta_0}^2 + \norm{h}^3\right) \label{eq:sum_forward_map_taylor_second}
    \end{align}
    as $(h, \delta_0) \to 0$, where $\delta_0 := g - g_0$ (where the remainder is uniform in $\lambda$ in the region where $|a_1+a_2|>a_\text{min}$).

    \item Since we would like an approximation that only depends on the representative atom $(g_0,a_0)$, 
    we will approximate $H_F(g)[h, h]$.
    Given some assumptions \hl{(such as $F$ has bounded third derivatives in a neighbourhood of $g_0$)}, 
    we have that \hl{(check!)}
    \[
    H_F(g)[h, h] = H_F(g_0)[h, h] + O(\norm{\delta_0}\norm{h}^2)
    \]
    as $(h, \delta_0) \to 0$ (where the remainder is uniform in $\lambda$ in the region where $|a_1+a_2|>a_\text{min}$).
    Combining this with \eqref{eq:sum_forward_map_taylor_second} gives
    \begin{align}
    \lambda F(g + \delta_1) + (1-\lambda)F(g + \delta_2)
    &= F(g_0) + J_F(g_0) \delta_0 + \frac{1}{2} \lambda(1-\lambda)H_F(g_0)[h, h] \nonumber\\
    &+ O\left(\norm{\delta_0}^2 + \norm{\delta_0}\norm{h}^2 + \norm{h}^3\right), \label{eq:sum_forward_map_taylor_third}
    \end{align}
    as $(h, \delta_0) \to 0$ (where the remainder is uniform in $\lambda$ in the region where $|a_1+a_2|>a_\text{min}$).

    \item Define $\eta := a - a_0$ to be the difference between the total amplitude and the amplitude of the representative atom.
    Then, from \eqref{eq:sum_forward_map_taylor_third} and \eqref{eq:sum_forward_map} we have that
    \begin{align}
    A(\theta_1) + A(\theta_2) 
    &= a\left[sF(c + \delta_1) + (1-s)F(c + \delta_2)\right] \nonumber\\
    &= (a_0 + \eta)\left[sF(c + \delta_1) + (1-s)F(c + \delta_2)\right] \nonumber\\
    &= (a_0 + \eta)\bigl[F(g_0) + J_{g_0} q + \frac{1}{2} s(1-s)H_{g_0}[h, h] \nonumber\\
    &+ O\left(\norm{q}^2 + \norm{q}\norm{h}^2 + \norm{h}^3\right)\bigr] \nonumber\\
    &= a_0F(g_0) + \eta F(g_0) + a_0 J_{g_0} q + \frac{a_0}{2} s(1-s)H_{g_0}[h, h] \nonumber\\
    &+ O\left(\abs{\eta}\norm{q} + \abs{\eta}\norm{h}^2 + \norm{q}^2 + \norm{q}\norm{h}^2 + \norm{h}^3\right) \nonumber\\
    &= A(\theta_0) + \eta F(g_0) + a_0 J_{g_0} q + \frac{a_0}{2} s(1-s)H_{g_0}[h, h] \nonumber\\
    &+ O\left(\abs{\eta}\norm{q} + \abs{\eta}\norm{h}^2 + \norm{q}^2 + \norm{q}\norm{h}^2 + \norm{h}^3\right) \nonumber
    \end{align}
    when $|a_1+a_2|>\epsilon$ . That is, 
    \begin{align}
    A(\theta_1) + A(\theta_2) - A(\theta_0)
    &= \eta F(g_0) + a_0 J_{g_0} q + \frac{a_0}{2} s(1-s)H_{g_0}[h, h] \nonumber\\
    &+ O\left(\abs{\eta}\norm{q} + \abs{\eta}\norm{h}^2 + \norm{q}^2 + \norm{q}\norm{h}^2 + \norm{h}^3\right) \label{eq:forward_map_approximation}
    \end{align}
    in some neighborhood around $\eta = a - a_0 = 0$. 

    \hl{where is it helpful that our parameterisation is as it is? so far it is just clear that making sure the deltas cancel gives the jacob 0, but not our particular choice}
\end{enumerate}


% -------- QUESTION 3 ---------
\pagebreak
\question*{}  
Suppose we have a representative atom $\theta_0=(x_{00},d_0,R_0,a_0)^{\mathsf T}\in\mathcal{X}_1$.
Provide a function that approximates the log posterior $\log p(\theta\mid y)$ on $\mathcal{X}_2$, 
where the observed data is the noise-free data $y=A(\theta_0)$ generated by this atom.

The function $f_2: \mathcal{X}_2 \to \R$ should satisfy the following properties:
\begin{enumerate}
    \item unimodal;
    \item exchangeable, i.e., $f_2(\theta_1, \theta_2) = f_2(\theta_2, \theta_1)$;
    \item a good approximation to the log posterior in the region of high probability around $\theta_0$;
    \item easy to compute defining quantities from $\theta_0$;
    \item easy to sample from;
    \item as simple as possible.
\end{enumerate}

Note that this function need only be specified up to an additive constant.
This should build on the function found in Question 2.

\textit{Solution.} 


% -------- QUESTION 4 ---------
\pagebreak
\question*{}  
Fix $n \geq 2$. 
Suppose we have a representative atom $\theta_0=(x_{00},d_0,R_0,a_0)^{\mathsf T}\in\mathcal{X}_1$.
Provide a function that approximates the log posterior $\log p(\theta\mid y)$ on $\mathcal{X}_n$, 
where the observed data is the noise-free data $y=A(\theta_0)$ generated by this atom. 

The function $f_n: \mathcal{X}_n \to \R$ should satisfy the following properties:
\begin{enumerate}
    \item unimodal;
    \item exchangeable;
    \item a good approximation to the log posterior in the region of high probability around $\theta_0$;
    \item easy to compute defining quantities from $\theta_0$;
    \item easy to sample from;
    \item as simple as possible.
\end{enumerate}

Note that this function need only be specified up to an additive constant, 
and this constant need not be common to all $n$.
This should be a generalization of the function found in Question 3.

\textit{Solution.} 


% -------- QUESTION 5 ---------
\pagebreak
\question*{}  
Suppose we have a representative atom $\theta_0=(x_{00},d_0,R_0,a_0)^{\mathsf T}\in\mathcal{X}_1$.
Rewrite the function found in Question 1 in the same form as 
the functions found in Question 4. 

If it is not possible to write the function found in Question 1 in the same form as 
the functions found in Question 4, explain why. Then, provide an alternative function 
that approximates the log posterior $\log p(\theta\mid y)$ on $\mathcal{X}_1$ 
that is in the same form as the functions found in Question 4.

\textit{Solution.} 


% -------- QUESTION 6 ---------
\pagebreak
\question*{}  
Suppose we have a representative atom $\theta_0=(x_{00},d_0,R_0,a_0)^{\mathsf T}\in\mathcal{X}_1$. 
Provide an approximation of the sequence of probabilities $\{p_{n}\}_{n\in\N}$ 
where 
\[
    p_n := P(\theta \in X_n \mid y) = \int_{\mathcal{X}_n} p(\theta\mid y) \rmd\theta, 
\]
where $y = A(\theta_0)$ is the noise-free data generated by $\theta_0$.

You may find it useful to use the functions found in Questions 4 and 5.
Note that you need only approximate the sequence up to a common multiplicative constant; 
that is, it suffices to approximate quantities proportional to $\{p_{n}\}_{n\in\N}$.

\textit{Solution.} 


% -------- QUESTION 7 ---------
\pagebreak
\question*{}  
Suppose we have a representative atom $\theta_0=(x_{00},d_0,R_0,a_0)^{\mathsf T}\in\mathcal{X}_1$. 
Provide an approximation of the probability $p$ 
where 
\[
    p := , 
\]
where $y = A(\theta_0)$ is the noise-free data generated by $\theta_0$.

\textit{Solution.} 


\end{document}
